\documentclass[11pt,oneside]{report}
\usepackage[a4paper,inner=26mm,outer=25mm,top=25mm,bottom=25mm,headheight=15pt,headsep=8mm]{geometry}
\usepackage[T1]{fontenc}
\usepackage[utf8]{inputenc}
\usepackage{amsmath,amsthm,mathtools}
\usepackage{newtxtext,newtxmath}
\usepackage{microtype,booktabs,longtable,array,enumitem,multicol}
\usepackage{xcolor,titlesec,fancyhdr}
\usepackage[most]{tcolorbox}
\usepackage{hyperref}
\usepackage[nameinlink,noabbrev]{cleveref}
\definecolor{ink}{HTML}{183443}
\definecolor{teal}{HTML}{136E71}
\definecolor{pale}{HTML}{EDF5F4}
\definecolor{graytext}{HTML}{53616B}
\hypersetup{hypertexnames=false,colorlinks=true,linkcolor=teal,citecolor=teal,urlcolor=teal,pdftitle={Pigeonhole Designs and Extension Elimination: A Research Compendium},pdfauthor={Research developed in conversation with ChatGPT},pdfsubject={Self-contained working lemmas, proofs, obstructions, and the route to AC0[p]-Frege lower bounds}}
\titleformat{\chapter}[display]{\normalfont\huge\bfseries\color{ink}}{\small\sffamily\bfseries CHAPTER \thechapter}{7pt}{}
\titlespacing*{\chapter}{0pt}{0pt}{22pt}
\titleformat{\section}{\Large\bfseries\color{ink}}{\thesection}{0.65em}{}
\titleformat{\subsection}{\large\bfseries\color{ink}}{\thesubsection}{0.65em}{}
\newcolumntype{L}[1]{>{\raggedright\arraybackslash}p{#1}}
\setlength{\parskip}{5pt}
\setlength{\parindent}{0pt}
\setlength{\emergencystretch}{3em}
\setlist{nosep,leftmargin=1.6em,topsep=4pt}
\pagestyle{fancy}\fancyhf{}
\fancyhead[L]{\small\sffamily\color{graytext}\nouppercase{\leftmark}}
\fancyfoot[L]{\footnotesize\sffamily\color{graytext}PHP / EXTENSION ELIMINATION}
\fancyfoot[R]{\small\thepage}
\renewcommand{\headrulewidth}{0.3pt}
\fancypagestyle{plain}{\fancyhf{}\fancyfoot[L]{\footnotesize\sffamily\color{graytext}PHP / EXTENSION ELIMINATION}\fancyfoot[R]{\small\thepage}\renewcommand{\headrulewidth}{0pt}}
\renewcommand{\chaptermark}[1]{\markboth{\thechapter\quad #1}{}}
\newtheoremstyle{research}{9pt}{5pt}{\itshape}{}{\bfseries}{.}{0.55em}{}
\theoremstyle{research}
\newtheorem{lemma}{Lemma}[chapter]
\newtheorem{theorem}[lemma]{Theorem}
\newtheorem{proposition}[lemma]{Proposition}
\newtheorem{corollary}[lemma]{Corollary}
\theoremstyle{definition}
\newtheorem{definition}[lemma]{Definition}
\newtheorem{example}[lemma]{Example}
\newtheorem{problem}[lemma]{Open target}
\newtheorem{imported}[lemma]{Imported input}
\theoremstyle{remark}
\newtheorem*{remark}{Scope}
\newtheorem*{editorial}{Editorial clarification}
\newtcolorbox{statusbox}{enhanced,breakable,colback=pale,colframe=teal,boxrule=0.5pt,arc=0pt,left=9pt,right=9pt,top=7pt,bottom=7pt}
\newcommand{\F}{\mathbb F}
\newcommand{\calF}{\mathcal F}
\newcommand{\calG}{\mathcal G}
\newcommand{\calB}{\mathcal B}
\newcommand{\calE}{\mathcal E}
\newcommand{\calR}{\mathcal R}
\newcommand{\BB}{\overline B}
\newcommand{\Ic}[1]{\mathcal I_{#1}}
\newcommand{\Cc}[1]{\mathcal C_{#1}}
\newcommand{\Jfld}{J_{\mathrm{fld}}}
\newcommand{\Fphp}{\mathcal F_n}
\newcommand{\PC}{\mathrm{PC}}
\newcommand{\NS}{\mathrm{NS}}
\newcommand{\AC}{AC^{0}[p]}
\DeclareMathOperator{\Span}{span}
\DeclareMathOperator{\rank}{rank}
\DeclareMathOperator{\im}{im}
\DeclareMathOperator{\rad}{rad}
\DeclareMathOperator{\Hilb}{Hilb}
\DeclareMathOperator{\NF}{NF}
\DeclareMathOperator{\supp}{supp}
\DeclareMathOperator{\diag}{diag}
\newcommand{\ledeg}[1]{_{\le #1}}
\newcommand{\pos}[1]{\left(#1\right)_{+}}
\newcommand{\chapterintro}[1]{\begin{statusbox}#1\end{statusbox}}
\begin{document}
\begin{titlepage}
\color{ink}
\vspace*{14mm}
{\small\sffamily\bfseries RESEARCH COMPENDIUM\par}
\vspace{11mm}
{\fontsize{31}{36}\selectfont\bfseries Pigeonhole Designs\par and Extension\par Elimination\par}
\vspace{8mm}
{\Large Working lemmas, proofs, and obstructions\par}
\vspace{5mm}
{\large Toward superpolynomial $\AC$-Frege lower bounds\par}
\vspace{15mm}
\begin{statusbox}
\textbf{What this document contains.} A consolidated mathematical record of the independent work developed in this conversation: the corrected low-degree algebra, generic annihilators, nonlinear liftings, polynomial-calculus elimination, nested and core/residual batching, exact decomposition optimization, and the remaining proof obligations.\par
\textbf{Current endpoint.} The static decomposition criterion is optimized, and an explicit admissible family defeats its affordable-budget target. The next missing theorem is refutation-sensitive elimination for PHP. No superpolynomial Frege lower bound is claimed here.
\end{statusbox}
\vfill
{\large Compiled 10 September 2026\par}
\vspace{4mm}
{\small\color{graytext}Source: the research conversation and its eleven computational supplements.\par Published inputs are attributed separately. Working proofs are not represented as independently refereed results.}
\end{titlepage}
\pagenumbering{roman}
\chapter*{How to read this record}
\addcontentsline{toc}{chapter}{How to read this record}
This is a mathematical consolidation, not a transcript and not a claim that the payoff has been achieved. It begins with the independent work undertaken after the initial uploaded research report. The report's older assertions are not silently imported as theorems; the corrections and changes of direction relevant to this conversation are recorded in Appendix~\ref{app:ledger}.

Each lemma supplies its assumptions and a proof of the claim made here. Repeated arguments have been factored into reusable preliminary lemmas; references between chapters are internal to this document. Results inherited from the literature are explicitly marked \emph{Imported input}. Their original proofs are not reproduced or claimed as our work. Statements involving the published PHP lower bound are conditional on the encoding match spelled out in Chapter~\ref{ch:route}.

\textbf{Status convention.} Unless marked imported or open, the results below are the \emph{working derivations obtained in the conversation}. Typesetting them as lemmas records their mathematical content; it does not certify novelty, independent verification, or a machine-checked general proof. A complete audit of the proof transformations remains an explicit obligation. The computational supplements establish only the finite checks described in their scope notes.

\textbf{Three separations matter throughout.} The non-Boolean graded ring $\BB$ is not the Boolean PHP ring. A design annihilating bounded-degree original-axiom multiples need not annihilate all bounded-degree polynomial-calculus consequences. Finally, generic linear restrictions are not the full leveled extension systems produced by a Frege simulation. None of these distinctions is erased by the notation.

\textbf{Degree convention.} An extension axiom retains its original degree in old and new variables. Specialization or Boolean reduction may lower the degree of its image, but does not retroactively enlarge the set of allowed multipliers in the original degree-$D$ design. Zero polynomials are omitted wherever a degree is used.

\textbf{Editorial changes.} Definitions, finite-dimensional duality, and degree-nonincreasing reductions are made explicit so the proofs can be followed without earlier messages. Conservative versus sharper degree bounds are identified. No missing global compatibility or simulation-transfer theorem has been supplied by assumption without being labeled.

\clearpage
\section*{The argument at a glance}
\begin{longtable}{@{}L{0.26\textwidth}L{0.67\textwidth}@{}}
\toprule
Branch & Endpoint and limitation\\\midrule
Graded algebra & Koszul vanishing and the initial Hilbert function hold in the proved range $2d-1\le n$. The originally proposed range $d<n$ is false.\\
Generic restrictions & A determinant functional survives $t$ generic restrictions when $t+2d-1\le n$. This concerns the non-Boolean ring only.\\
Exact liftings & Reweighting, factor packing, shared conditioning, base-aware substitutions, and supported moment corrections give explicit constructions with stated budgets.\\
PC elimination & Blocks can be eliminated at additive degree cost. Equal spans, nested spans, and nested cores with low-rank residuals can be batched more cheaply.\\
Decomposition & For affine data the current cost criterion is exactly optimizable. There are polynomial-size admissible families whose optimum still exceeds $n/2$.\\
Remaining theorem & Use the actual translated Nullstellensatz refutation and PHP structure to obtain a base PC refutation below the available linear-degree threshold.\\\bottomrule
\end{longtable}
\section*{Notation quick reference}
\begin{tabular}{@{}L{0.22\textwidth}L{0.71\textwidth}@{}}
$\BB$ & Non-Boolean column-exclusive ring with equal row sums.\\[4pt]
$\mathcal F_n$ & Boolean linear-row PHP axioms, without row exclusion.\\[4pt]
$\mathcal I_D$ & Span of original-axiom multiples of total degree at most $D$.\\[4pt]
$\mathcal C_D$ & Polynomial-calculus consequences derivable within degree $D$.\\[4pt]
$h,\ \delta,\ D$ & Extension accuracy, maximum input degree, and target proof degree.\\[4pt]
$B,\ T$ & Required base degree and polynomial-substitution degree inflation.\\
\end{tabular}

\textbf{Reading paths.} Chapters 1--2 contain the graded-algebra branch. Chapters 3--6 develop explicit functionals and their limitations. Chapters 7--9 give the current proof-elimination route and its decomposition obstruction. Chapter 10 states the remaining goal and risk-ordered proof obligations; Appendix C provides a compact budget reference.
\tableofcontents
\clearpage
\pagenumbering{arabic}
\chapter{Objects, degrees, and reusable tools}\label{ch:foundations}
\chapterintro{This chapter fixes the exact rings and proof systems, then proves the elementary tools used by every later construction. The distinction between a truncated consequence space and an unrestricted ideal is essential.}

\section{The base systems}
Let $k$ be a field. For $m$ rows and $n$ columns define
\[
 S_{m,n}=k[x_{ij}:i\in[m],j\in[n]]/
 (x_{ij}x_{i'j}:i\ne i').
\]
Variables have degree one. Rows may repeat, and powers of variables are allowed. Only different rows in the same column are excluded. Put $\rho_i=\sum_jx_{ij}$. For $m=n+1$, write
\[
 S=S_{n+1,n},\qquad
 \BB=S/(\rho_i-\rho_1:2\le i\le n+1),\qquad z=\rho_1\in\BB.
\]

For fixed prime $p$, our Boolean linear-row base system over $k=\F_p$ is
\begin{equation}\label{eq:base}
 \Fphp=\{\rho_i-1\}_i\cup
 \{x_{ij}x_{i'j}:i\ne i'\}\cup\{x_{ij}^2-x_{ij}\}_{i,j}.
\end{equation}
No same-row exclusion equations are assumed. A Boolean row with sum one modulo $p$ need not have exactly one occupied cell. The homogeneous Boolean ring relevant to \eqref{eq:base} is
\[
 A_n=k[x,z]/(\rho_i-z,\ x_{ij}x_{i'j},\ x_{ij}^2-zx_{ij})
     \cong\BB/(x_{ij}^2-zx_{ij}).
\]
Results in $\BB$ do not automatically descend to $A_n$.

\begin{lemma}[Axiom deletion]\label{lem:deletion}
If $\calF\subseteq\calF^+$, every degree-$D$ design for $\calF^+$ is one for $\calF$. Every PC refutation of $\calF$ is also a refutation of $\calF^+$ with the same degree. Consequently a degree lower bound for the stronger axiom system transfers to the weaker one.
\end{lemma}
\begin{proof}
The design has fewer equations to annihilate after deletion. A proof using only the smaller axiom list is still a valid proof from the larger list. Taking the contrapositive gives the lower-bound assertion.
\end{proof}

\section{Designs and polynomial calculus}
For a finite set $\calF\subseteq k[y]$, define
\[
 \Ic{D}(\calF)=\Span_k\{qf:f\in\calF,\ \deg(qf)\le D\}.
\]
A degree-$D$ \emph{design} is a linear map $\lambda:k[y]_{\le D}\to k$ with
\[
 \lambda(1)=1,\qquad \lambda(\Ic{D}(\calF))=0.
\]
A degree-$D$ Nullstellensatz refutation is an identity
$1=\sum_fq_ff$ with every $\deg(q_ff)\le D$.

Polynomial calculus (PC) starts from the axioms, and permits field-linear combinations and multiplication of a previously derived polynomial by one variable. Its degree is the largest total degree of any line, after collecting like terms in the ordinary polynomial ring. Let $\Cc{D}(\calF)$ be the vector space of polynomials having PC derivations within degree $D$.

\begin{lemma}[Finite-dimensional separation]\label{lem:duality}
A degree-$D$ design for $\calF$ exists exactly when $1\notin\Ic{D}(\calF)$. A normalized functional annihilating $\Cc{D}(\calF)$ exists exactly when there is no degree-$D$ PC refutation. Moreover,
\[
 \Ic{D}(\calF)\subseteq\Cc{D}(\calF).
\]
\end{lemma}
\begin{proof}
A normalized functional vanishing on a subspace cannot exist if that subspace contains $1$. Conversely, extend a basis of the subspace by $1$ and then to a basis of $k[y]_{\le D}$; assign zero to the original basis, one to $1$, and arbitrary values to the remaining basis elements. For the inclusion, derive each monomial multiple of an original axiom by repeated variable multiplication and combine the results. All intermediate degrees are bounded by the final product degree. The set of PC-derivable polynomials of bounded degree is closed under linear combinations by concatenating derivations.
\end{proof}

\begin{lemma}[Reuse versus flattening]\label{lem:reuse}
If $f\in\Cc{c}(\calF)$ and $q$ is a polynomial, then
\[
 qf\in\Cc{\max\{c,\deg q+\deg f\}}(\calF).
\]
In contrast, the immediate consequence for a truncated original-axiom representation is only
\[
 q\Ic{c}(\calF)\subseteq\Ic{c+\deg q}(\calF).
\]
\end{lemma}
\begin{proof}
For PC, keep the derivation of $f$, multiply its \emph{final line} by each monomial of $q$, and form their linear combination. No old derivation line is multiplied by $q$. In a flattened certificate for $f$, multiplying each certificate summand by $q$ can raise the degree of every such summand by $\deg q$. These are different operations and yield the two bounds.
\end{proof}

\begin{lemma}[Polynomial substitutions in PC]\label{lem:substitution}
Let $\Phi$ map variables to polynomials of degree at most $T\ge1$. A PC derivation of degree $D$ maps to a PC derivation of degree at most $TD$ from the substituted axioms. If every line is also multiplied by a fixed polynomial $W$, the ceiling becomes $TD+\deg W$, provided the weighted substituted axiom lines have derivations within that ceiling.
\end{lemma}
\begin{proof}
Linear combinations commute with substitution. A multiplication step $P\mapsto yP$ maps to $\Phi(P)\mapsto\Phi(y)\Phi(P)$. Expand $\Phi(y)$ into monomials and use repeated variable multiplication followed by linear combination. Since the original predecessor has degree at most $D-1$, all these intermediate products have degree at most $T(D-1)+T$. The same argument with the fixed factor $W$ proves the weighted version. Zero predecessors cause no difficulty.
\end{proof}

\section{Finite-field identities and bounded-degree reductions}
Every old variable is assigned a domain, either $\{0,1\}$ or $\F_p$. Let $\Jfld$ be generated by $y^2-y$ or $y^p-y$, respectively.

\begin{lemma}[Degree-nonincreasing field reduction]\label{lem:fieldreduction}
A polynomial $P$ that vanishes on the product of the specified old-variable domains belongs to $\Jfld$. If $\deg P\le B$, it has an ideal representation by the defining univariate equations in which every summand has degree at most $B$. In particular, a degree-$B$ design annihilates $P$. For every $g\in\F_p[y]$,
\[
 g^p-g\in\Jfld.
\]
\end{lemma}
\begin{proof}
Divide successively by the monic univariate generators. Each replacement lowers a power of one variable and never increases total degree; record each subtracted multiple. The remainder has individual degree below the size of the corresponding domain. A polynomial with those individual-degree bounds and zero values on the full product grid is zero: apply univariate interpolation in one coordinate and induct on the number of coordinates. Thus the remainder vanishes and the recorded representation has the stated degree bound. Finally, every old assignment makes $g$ an element of $\F_p$, so $g^p-g$ vanishes on the domain.
\end{proof}

\begin{lemma}[Finite-field selectors]\label{lem:selectors}
For $\alpha\in\F_p$ put $\chi_\alpha(t)=1-(t-\alpha)^{p-1}$. Then
\begin{align}
 \sum_{\alpha\in\F_p}\chi_\alpha(t)&=1,\label{eq:selectorpartition}\\
 t&=\sum_{\alpha\ne0}\alpha\chi_\alpha(t),\label{eq:selectorlinear}\\
 \chi_\alpha(t)(1-\alpha^{-1}t)&=\alpha^{-1}(t^p-t)
       \quad(\alpha\ne0),\label{eq:selectorproduct}\\
 1-t^{p-1}&=\prod_{\alpha\ne0}(1-\alpha^{-1}t).
 \label{eq:factorzero}
\end{align}
All are ordinary polynomial identities over $\F_p$.
\end{lemma}
\begin{proof}
The first two identities have degree at most $p-1$ and agree on $\F_p$: the $\chi_\alpha$ are its point indicators. This also covers $p=2$. For the third, write $1-\alpha^{-1}t=-\alpha^{-1}(t-\alpha)$ and use $(t-\alpha)^p=t^p-\alpha$. Both sides of the fourth have degree $p-1$, the same nonzero roots, and constant term one.
\end{proof}

\begin{lemma}[Column-exclusive normal form]\label{lem:columnnf}
Let
\[
 J=(x_{ij}^2-x_{ij},\ x_{ij}x_{i'j}:i\ne i').
\]
Its quotient is the algebra of functions on assignments where every column is either empty or occupied by one row. Every polynomial in $J$ of degree at most $B$ has a representation by its displayed generators within degree $B$.
\end{lemma}
\begin{proof}
Reduce powers by $x_{ij}^2\mapsto x_{ij}$ and different-row products in a column to zero. The surviving monomials choose at most one variable per column. For a single column, $1,x_{1j},\ldots,x_{mj}$ are independent functions on its $m+1$ states; tensoring proves uniqueness on all columns. The reductions never increase degree and record the desired certificate.
\end{proof}

\section{Extension data and the actual quantifiers}
An extension block $a$, of accuracy $h_a\ge1$, has old-variable inputs $g_{a,1},\ldots,g_{a,k_a}$, fresh variables $r_{a,u,j}$, and all companions
\begin{equation}\label{eq:ens}
 E_{a,i}=g_{a,i}\prod_{u=1}^{h_a}
 \left(1-\sum_jr_{a,u,j}g_{a,j}\right).
\end{equation}
Within one level, new variables are disjoint across blocks. Inputs use only original variables and variables from earlier levels. The field axioms $\calR=\{r^p-r\}$ are included. This definition and level convention follow \cite[Definition 2.1]{Krajicek}.

With zero inputs omitted, put
\[
 \delta_{a,i}=\deg g_{a,i},\quad \delta_a=\max_i\delta_{a,i},\quad
 e_{a,i}=\delta_{a,i}+h_a(\delta_a+1).
\]
Fresh coefficient variables make this the original degree of a nonzero companion. A companion is \emph{active at $D$} if $e_{a,i}\le D$.

\begin{lemma}[Scalar lifting is sufficient, not necessary]\label{lem:scalar-sufficient}
Suppose $\beta\in\F_p^{|r|}$ and a normalized functional $\lambda$ on old polynomials annihilates the images of every allowed degree-$D$ axiom multiple under $r=\beta$. Then
\[
 \Lambda(P)=\lambda(P(y,\beta))
\]
is a degree-$D$ design for the original joint-variable system. An arbitrary joint design need not have this form.
\end{lemma}
\begin{proof}
Constant substitution preserves degree, commutes with multiplication, and makes $r^p-r$ zero. The assumed image equations are exactly the design equations. No assertion about representing an arbitrary linear functional by one evaluation follows from these facts.
\end{proof}

The desired existential statement is $\forall\calE\ \exists\Lambda_{\calE}$, for the extension data actually required by the simulation. The stronger version $\forall\calE\ \exists\beta\ \exists\lambda$ is one sufficient route. Universal success of all scalar choices is false: for the input tuple $(1)$, the companion is $\prod_u(1-r_u)$, which becomes $1$ at $r=0$ and $0$ when $r_1=1$.

\begin{lemma}[Homogenized separation]\label{lem:homogenization}
For the base system \eqref{eq:base} and additional old-variable polynomials $f$, homogenize each nonzero $f$ to its actual degree. There is a degree-$D$ design precisely when
\[
 z^D\notin(f^{\mathrm h})A_n.
\]
Membership is tested in homogeneous degree $D$.
\end{lemma}
\begin{proof}
Homogenizing a bounded-degree certificate $1=\sum q_ff$ to degree $D$ produces $z^D$ in the displayed ideal, including the homogenized base equations. Conversely, dehomogenize any homogeneous degree-$D$ representation by setting $z=1$; every certificate summand has degree at most $D$. Apply \cref{lem:duality}. This equivalence concerns $A_n$, not the larger non-Boolean ring $\BB$.
\end{proof}

\chapter{The graded branch: corrected vanishing and generic survival}\label{ch:graded}
\chapterintro{This branch is a result about $\BB$, before Boolean equations are imposed. It supplies the corrected form of item (a), a weak restriction count, and an explicit generic annihilator. It does not by itself solve the extension-system problem.}

\section{The original range fails}
For linear forms $\rho_1,\ldots,\rho_m$ in $S_{m,n}$, grade the Koszul complex by
\[
 K_i(\rho;S)_j=S_{j-i}\otimes_k\bigwedge^i k^m.
\]
Give $x_{ab}$ and the $a$th exterior generator row degree $\varepsilon_a$. Its differential preserves this finer grading.

\begin{lemma}[A characteristic-free counterexample]\label{lem:koszul-counter}
For $n=4,m=5$, the row multidegree $\alpha=(1,1,1,0,0)$ satisfies
\[
 H_1(\rho_1,\ldots,\rho_5;S_{5,4})_\alpha\ne0.
\]
Thus the statement $H_i(\rho;S)_j=0$ for every $i\ge1$ and $j<n$ is false over every field.
\end{lemma}
\begin{proof}
The strand has dimensions
\[
 0\longrightarrow k\xrightarrow{\partial_3}k^{12}
 \xrightarrow{\partial_2}k^{36}
 \xrightarrow{\partial_1}k^{24}\longrightarrow0.
\]
The counts are respectively $1$, $3\cdot4$, $3\cdot4\cdot3$, and $4\cdot3\cdot2$: active rows remaining in the coefficient monomial must use distinct columns. The image of the unique top exterior basis element is
$\rho_1e_2\wedge e_3-\rho_2e_1\wedge e_3+\rho_3e_1\wedge e_2$, which is nonzero. Hence $\rank\partial_3=1$ and $\rank\partial_2\le11$. But $\dim\ker\partial_1\ge36-24=12$, giving $\dim H_1\ge1$. The internal degree is $3<4$.
\end{proof}

\section{A multigraded vanishing lemma}
\begin{lemma}[Vanishing with active-row control]\label{lem:multigraded}
Let $S_{r,N}$ have $r$ rows and $N$ columns, let $\sigma_a$ be its row sums, and suppose $\alpha_a>0$ for all $a$. Put $q=\sum_a\alpha_a$. For $i\ge1$,
\[
 \boxed{N\ge q+r-i\quad\Longrightarrow\quad
 H_i(\sigma;S_{r,N})_\alpha=0.}
\]
\end{lemma}
\begin{proof}
Write $T=S_{r,N-1}$ and
$C_r=k[y_1,\ldots,y_r]/(y_ay_b:a\ne b)$ for the final column. There is an exact sequence
\begin{equation}\label{eq:axis-exact}
 0\longrightarrow C_r\longrightarrow\bigoplus_{a=1}^r k[y_a]
 \longrightarrow k^{r-1}\longrightarrow0.
\end{equation}
The first map restricts to the coordinate axes; its image consists of tuples with equal constant terms. The last map takes differences of those constants.

Tensor with $T$. Let $\sigma'_b$ be a row sum on the first $N-1$ columns. On the middle summand $T[y_a]$, the full row sum acts as $\sigma'_b+\delta_{ab}y_a$; on the final $T^{r-1}$ it acts as $\sigma'_b$. Thus \eqref{eq:axis-exact} induces an exact sequence of Koszul complexes.

On $T[y_a]$, the element $\sigma'_a+y_a$ is monic in $y_a$ and is a nonzerodivisor. Its two-term Koszul complex resolves the quotient $T[y_a]/(\sigma'_a+y_a)\cong T$. Tensoring the other, free Koszul factors and taking homology identifies this middle summand with the Koszul complex on $\sigma'_b$, $b\ne a$, over $T$.

Abbreviate $H_i^{r,N}(\alpha)=H_i(\sigma;S_{r,N})_\alpha$. The resulting long exact sequence contains
\begin{equation}\label{eq:rowinduction}
 (H_{i+1}^{r,N-1}(\alpha))^{r-1}
 \longrightarrow H_i^{r,N}(\alpha)
 \longrightarrow\bigoplus_a H_i(\sigma'_b:b\ne a;T)_\alpha.
\end{equation}
Fix a row $a$ in the right-hand term. Its row sum has been omitted, so the differential does not change the row-$a$ part $\mu$ of a coefficient monomial. This monomial has degree $\alpha_a$ and occupies $s(\mu)\le\alpha_a$ columns. All other rows must avoid those columns. Consequently
\[
 H_i(\sigma'_b:b\ne a;T)_\alpha
 \cong\bigoplus_{\deg\mu=\alpha_a}
 H_i^{r-1,N-1-s(\mu)}(\alpha_{\widehat a}).
\]
Repeated powers and repeated use of a row are allowed in this decomposition.

Induct on $r$, and, at fixed $r$, downwards on $i$. For one row the ring is a polynomial ring and its nonzero row sum is a nonzerodivisor; homology above index one is automatically zero. Zero-row strands also have no positive homology. If $N\ge q+r-i$, the left term of \eqref{eq:rowinduction} vanishes by downward induction since $N-1\ge q+r-(i+1)$. Each right summand vanishes because
\[
 N-1-s(\mu)\ge(q-\alpha_a)+(r-1)-i.
\]
Exactness forces the middle term to vanish.
\end{proof}

\begin{corollary}[Corrected item (a): Koszul range]\label{cor:koszul}
For any number of rows and $n$ columns,
\[
 H_i(\rho;S_{m,n})_j=0
 \quad\text{if }i\ge1\text{ and }2j-1\le n.
\]
\end{corollary}
\begin{proof}
A row multidegree of total degree $j$ has at most $r\le j$ positive coordinates. Rows with zero coordinate cannot occur in a coefficient monomial or an exterior generator, so the strand is the $r$-row strand of \cref{lem:multigraded}. Since $j+r-i\le2j-1$, that lemma applies. Sum over all row multidegrees.
\end{proof}

\begin{lemma}[Initial Hilbert function and $z$-injectivity]\label{lem:hilbert}
For $m=n+1$ and $0\le d\le\lfloor(n+1)/2\rfloor$,
\[
 \boxed{\dim_k\BB_d=\binom nd n^d.}
\]
For $1\le d\le\lfloor(n+1)/2\rfloor$, multiplication by $z$ maps $\BB_{d-1}$ injectively into $\BB_d$. In particular, $z^d\ne0$ in this range. No characteristic restriction is needed.
\end{lemma}
\begin{proof}
Put $\ell_i=\rho_i-\rho_1$ for $i\ge2$. An invertible change of generators identifies $K(\rho;S)$ with $K(\ell,z;S)$. Appending $z$ is the mapping cone of multiplication by $z$ on $K(\ell;S)$. Its long exact sequence contains
\[
 H_i(\ell;S)_{d-1}\xrightarrow{z}H_i(\ell;S)_d
 \longrightarrow H_i(\rho;S)_d
 \longrightarrow H_{i-1}(\ell;S)_{d-1}.
\]
Starting at internal degree zero, \cref{cor:koszul} and induction on $d$ give $H_i(\ell;S)_d=0$ for every $i\ge1$ in the stated range.

One column has Hilbert series $1+mt/(1-t)=(1+(m-1)t)/(1-t)$. Independent columns tensor, hence
\[
 \Hilb_S(t)=\frac{(1+nt)^n}{(1-t)^n}.
\]
The Euler series for the Koszul complex on the $n$ row differences is $(1-t)^n\Hilb_S(t)=(1+nt)^n$. Positive homology vanishes in the range under consideration, so its coefficients there equal those of $H_0=\BB$. This gives the Hilbert function. Finally the long exact sequence has
$H_1(\rho;S)_d\to\BB_{d-1}\xrightarrow{z}\BB_d$; its first term is zero. Since $\BB_0=k$, repeated injectivity proves $z^d\ne0$.
\end{proof}

\begin{lemma}[The weak restriction count is automatic]\label{lem:weakcount}
Let $L_1,\ldots,L_t$ be arbitrary linear forms in $\BB$. In the range of \cref{lem:hilbert},
\[
 \dim(\BB/(L_1,\ldots,L_t))_d
 \ge \binom n{d-1}n^{d-1}
 \left(\frac{n(n-d+1)}d-t\right).
\]
In particular this dimension is positive for $t<n$ and $1\le d\le\lfloor(n+1)/2\rfloor$.
\end{lemma}
\begin{proof}
The generated degree-$d$ subspace is the image of
$\BB_{d-1}^{\oplus t}\to\BB_d$, $(h_s)\mapsto\sum_sL_sh_s$.
Its rank is at most $t\dim\BB_{d-1}$. Subtract and use the ratio
$\dim\BB_d/\dim\BB_{d-1}=n(n-d+1)/d\ge n$.
\end{proof}

\begin{lemma}[Affine feasibility and its quadratic form]\label{lem:affinefeasibility}
For linear forms $L_1,\ldots,L_t$, the following are equivalent:
\[
 z^d\notin(L_1,\ldots,L_t)\BB,\qquad
 \exists\lambda:\BB_d\to k:\ 
 \lambda(z^d)=1,\quad \lambda(L_s\BB_{d-1})=0.
\]
At $d=2$, this is equivalent to a symmetric bilinear form $B$ on $\BB_1$ such that
\[
 B(x_{ij},x_{i'j})=0\ (i\ne i'),\qquad
 L_s\in\rad B,\qquad B(z,z)=1.
\]
\end{lemma}
\begin{proof}
The first equivalence is finite-dimensional separation of the degree-$d$ relation image from $z^d$. At degree two, a functional on the symmetric square of $\BB_1$ is the same as a symmetric bilinear form via $B(v,w)=\lambda(vw)$, including in characteristic two. It descends to $\BB_2$ exactly when it kills the column-exclusion quadratics. Killing the $L_s$ multiples is exactly the radical condition, and normalization is $B(z,z)=1$.
\end{proof}

\section{The generic annihilating functional}
\begin{lemma}[A special maximal-rank restriction family]\label{lem:special-columns}
Suppose $t+2d-1\le n$. Let $\sigma_s=\sum_i x_{is}$ be the sums of the first $t$ columns. Then
\[
 H_i(\ell_2,\ldots,\ell_{n+1},\sigma_1,\ldots,\sigma_t;S)_j=0
 \quad(i\ge1,\ j\le d),
\]
and $z^d$ survives the quotient by these linear forms.
\end{lemma}
\begin{proof}
In one column $C_m$, multiplication by $\sigma=\sum_i y_i$ is injective: write an element uniquely as a constant plus $\sum_i f_i(y_i)$ with $f_i(0)=0$ and compare the separate-axis monomials in its product with $\sigma$. Moreover
\[
 C_m/(\sigma)\cong k\oplus V,\qquad \dim V=m-1,\qquad V^2=0,
\]
because $\sigma y_i=0$ forces $y_i^2=0$. The different $\sigma_s$ are in separate tensor factors and form a regular sequence. The quotient is
\[
 T=S_{m,n-t}\otimes Q,\qquad Q=(k\oplus V)^{\otimes t}.
\]
In $T$, a full row sum is $\rho'_i+u_i$ with $\rho'_i$ on the remaining columns and $u_i\in Q_1$. Filter its Koszul complex by $Q$-degree. The part preserving that degree is the row-sum complex on $S_{m,n-t}$, tensored with a graded piece of $Q$. By \cref{cor:koszul} it has zero positive homology through degree $d$.

For completeness, a positive-index cycle can be killed by taking its lowest $Q$-degree part, writing that part as a boundary in the associated graded complex, subtracting the corresponding full boundary, and repeating. At fixed total degree the filtration is finite. This proves positive homology vanishing for the full row sums in $T$, equivalently for $(\rho,\sigma)$ in $S$.

Replace $\rho$ by $(\ell,z)$ and use the mapping-cone argument from \cref{lem:hilbert}. It proves both vanishing for $(\ell,\sigma)$ and injectivity of $z$ in its quotient through degree $d$. The degree-zero quotient is $k$, so $z^d$ survives.
\end{proof}

\begin{theorem}[Generic survival with an explicit determinant]\label{thm:generic}
For $t+2d-1\le n$, there is a nonzero polynomial $\Delta$ in the coefficients of $t$ linear forms $L_s$ such that, whenever $\Delta\ne0$, a functional $\lambda:\BB_d\to k$ satisfies
\[
 \lambda(z^d)=1,\qquad
 \lambda(L_su)=0\quad(s\le t,\ u\in\BB_{d-1}).
\]
The nonvanishing locus contains a specified $k$-rational point: the first $t$ column sums.
\end{theorem}
\begin{proof}
In the monomial basis of $S_d$, let $A(C)$ have columns $\ell_i\mu$ and $L_s(C)\mu$, for every degree-$(d-1)$ basis monomial $\mu$. Let $v$ be the coordinate vector of $\rho_1^d$. At the column-sum array $C_0$, put $r=\rank A(C_0)$. By \cref{lem:special-columns}, $v\notin\im A(C_0)$.

Let $B(C)$ be the preceding Koszul differential into the domain of $A(C)$. The chain identity gives $A(C)B(C)=0$. At $C_0$, exactness at this domain gives
\[
 \rank A(C_0)+\rank B(C_0)=\dim K_{1,d}.
\]
Over $k(C)$ both ranks are at least their specialized values, because a minor nonzero at $C_0$ is a nonzero polynomial. Their sum can never exceed $\dim K_{1,d}$. Neither can increase, so the generic rank of $A$ is $r$, and every specialization has rank at most $r$.

Choose $r$ columns $J$ of $A(C_0)$ that are independent, then $r+1$ row positions $R$ such that
\[
 \Delta(C)=\det[\,A(C)_{R,J}\mid v_R\,]
\]
is nonzero at $C_0$. For a degree-$d$ polynomial $f$ define
\begin{equation}\label{eq:detfunctional}
 \lambda_C([f])=
 \frac{\det[\,A(C)_{R,J}\mid [f]_R\,]}{\Delta(C)}.
\end{equation}
It is linear in $f$ and takes value one on $v$. When $\Delta(C)\ne0$, the selected $r$ columns span the entire relation image, since its rank is at most $r$. Substituting any relation into the last column makes the numerator zero. Thus the functional descends through all the required relations.
\end{proof}

\begin{corollary}[Affine design from the homogeneous witness]\label{cor:affinegeneric}
For $f=\sum_{e=0}^df_e$ with homogeneous $f_e$, the formula
\[
 \Lambda_C(f)=\lambda_C\left(\sum_{e=0}^dz^{d-e}f_e\right)
\]
is normalized and annihilates all degree-at-most-$d$ multiples of column exclusion, $\rho_i-1$, and the chosen $L_s=0$ equations.
\end{corollary}
\begin{proof}
Homogenize to degree $d$. A row-normalization multiple contains $\rho_i-z$, zero in $\BB$; a restriction multiple lies in the relation image of \cref{thm:generic}; column products are already zero. The constant one homogenizes to $z^d$.
\end{proof}

\begin{lemma}[Why genericity and dimension are insufficient]\label{lem:linearobstructions}
One unrestricted form can kill the target: $L_1=z$ implies $z^d=0$ in the quotient. Even excluding $z$ from the linear span of the restrictions does not suffice. For $n\ge2$, take
\[
 L_1=z-x_{11},\qquad L_2=z-x_{21}.
\]
Their span does not contain $z$, but $z^2\in(L_1,L_2)\BB$.
\end{lemma}
\begin{proof}
The first assertion is $z^d=L_1z^{d-1}$. For the second,
\[
 zL_1+x_{11}L_2=z^2-x_{11}x_{21}=z^2.
\]
To check the linear-span claim, compare coefficients modulo the row-difference subspace: within each row a combination of row sums has equal coefficients in every column. A combination $a(z-x_{11})+b(z-x_{21})=z$ would force the column-1 coefficient deviations in the first two rows to be zero, hence $a=b=0$, after which $z=0$ would be required in degree one. But \cref{lem:hilbert} gives $z\ne0$.
\end{proof}

\begin{remark}
A nonzero $\Delta$ over $\F_p$ need not be nonzero on every constrained family of $\F_p$-points. The known good point $C_0$ need not lie in the coefficient family arising from an extension construction. The surviving quotient can be large while its distinguished element has already died. These facts are why the later chapters change the object being constructed.
\end{remark}

\input{03_reweighting.tex}
\chapter{Factor packing and pointwise barriers}\label{ch:packing}
\chapterintro{Polynomial substitutions can use all $h$ extension factors in parallel. This removes block count from one budget, but not the worst-case rank of the input tuple. Shared weights and substitutions can be combined; the barriers here delimit those pointwise mechanisms.}

\section{Exact factor packing}
\begin{lemma}[Packed zero-test coefficients]\label{lem:packedcoeff}
Let $g_1,\ldots,g_k$ span a polynomial space of dimension $R$, and choose a basis $f_1,\ldots,f_R$ from the tuple, of degree at most $\delta$. For accuracy $h$, there are polynomial coefficients $\beta_{u,j}(y)$ such that
\[
 \prod_{u=1}^h\left(1-\sum_j\beta_{u,j}g_j\right)
 =Z:=\prod_{\nu=1}^R(1-f_\nu^{p-1}),
\]
with
\[
 \boxed{\deg\beta_{u,j}\le
 \delta\left(\left\lceil\frac{(p-1)R}{h}\right\rceil-1\right)_+.}
\]
Moreover $g_iZ\in\Jfld$ for every $i$.
\end{lemma}
\begin{proof}
By \eqref{eq:factorzero}, $Z$ is a product of $(p-1)R$ factors $1-\alpha^{-1}f_\nu$. Partition them into $h$ groups of at most $Q=\lceil(p-1)R/h\rceil$ factors. For a group ordered as $(\alpha_\ell,\nu_\ell)$, define the coefficient of $f_\nu$ by
\[
 \beta_{u,\nu}=
 \sum_{\ell:\nu_\ell=\nu}\alpha_\ell^{-1}
 \prod_{v<\ell}(1-\alpha_v^{-1}f_{\nu_v}).
\]
The identity
$1-\prod_{\ell=1}^b(1-a_\ell)=\sum_{\ell=1}^b a_\ell\prod_{v<\ell}(1-a_v)$
proves the group formula, and the largest coefficient uses at most $Q-1$ inputs. An empty group uses zero coefficients. Convert the chosen basis back to original input coordinates by constants. Finally,
$f_\nu Z=(f_\nu-f_\nu^p)\prod_{\mu\ne\nu}(1-f_\mu^{p-1})\in\Jfld$;
linearity gives the assertion for every $g_i$.
\end{proof}

\begin{lemma}[A level costs a maximum, not a sum]\label{lem:packedlift}
For blocks $a$ let $R_a$ and $\delta_a$ be their input ranks and degree bounds, and put
\[
 T_a=\max\left\{1,\delta_a
 \left(\left\lceil\frac{(p-1)R_a}{h_a}\right\rceil-1\right)_+\right\},
 \qquad T=\max_aT_a.
\]
A degree-$TD$ old design lifts to a degree-$D$ design for the entire level via
\[
 \Lambda(P)=\lambda(P(y,\beta(y))).
\]
It preserves all old moments through degree $D$.
\end{lemma}
\begin{proof}
Use \cref{lem:packedcoeff} in every block simultaneously. All substituted variables depend only on old variables, so no within-level nesting occurs and a degree-$D$ polynomial has image degree at most $TD$. Extension images belong to $\Jfld$, and a new field equation has image $\beta^p-\beta\in\Jfld$. Degree-bounded reduction handles their multiples. Old axioms remain old-axiom multiples. Substitution fixes old polynomials and one.
\end{proof}

\begin{corollary}[Variants and level composition]\label{cor:packingvariants}
If the basis inputs are Boolean-valued modulo $\Jfld$, replace $(p-1)R_a$ by $R_a$. Only active companions need be spanned when constructing a degree-$D$ design. Across levels the sufficient degree is $D\prod_jT_j$. For affine inputs over $\F_2$, rank $R_a\le2h_a$ gives $T_a=1$, regardless of the number of blocks.
\end{corollary}
\begin{proof}
For a Boolean-valued $f$, use $1-f$ instead of $1-f^{p-1}$. Inactive companions have no permitted multipliers, so a zero test for the active inputs suffices. Compose degree inflation one level at a time, retaining the original axiom degrees. For affine inputs and $R_a\le2h_a$, $\lceil R_a/h_a\rceil-1\le1$.
\end{proof}

\begin{lemma}[Packing plus repair]\label{lem:hybrid}
Fix $\tau\ge1$. Pack all blocks with $T_a\le\tau$, and reserve the others for scalar repair. Define
\[
 c_a^{(\tau)}(D)=\max\left(\{0\}\cup
 \{\tau(D-e_{a,i})+(p-1)\delta_{a,i}:e_{a,i}\le D\}\right).
\]
A sufficient old-design degree is
\[
 \boxed{B_{\rm hybrid}(D)=
 \min_{\tau\ge1}\left[\tau D+\sum_{a:T_a>\tau}c_a^{(\tau)}(D)\right].}
\]
In particular it is no worse than either $D+\sum_ac_a(D)$ or $D\max_aT_a$.
\end{lemma}
\begin{proof}
After packed substitutions an allowed cofactor of $E_{a,i}$ has old-variable degree at most $\tau(D-e_{a,i})$. In the repair procedure of \cref{lem:levelrepair}, test all moments $\lambda(Wqg_{a,i})$ up to that enlarged degree. Each repair costs at most $c_a^{(\tau)}(D)$. Packed blocks already vanish modulo $\Jfld$, so reweighting cannot spoil them. The final weighted substituted polynomial has degree at most $\deg W+\tau D$. At $\tau=1$ this only deletes costs from the old repair bound; at $\tau=\max T_a$ no repaired blocks remain. Between consecutive $T_a$ values the displayed cost is nondecreasing, so only those thresholds and $1$ need be tested.
\end{proof}

\begin{lemma}[Packing plus shared conditioning]\label{lem:sharedhybrid}
For one level, suppose one group is handled by substitutions of degree at most $T$, and all remaining blocks are determined by common features with weight cost $\kappa$. Then $B=TD+\kappa$ suffices.
\end{lemma}
\begin{proof}
Choose one nonzero shared weight $W$ and scalar values for its blocks, and use the polynomial substitutions on the first group. Define
$\Lambda(P)=\lambda(WP(y,\beta_{\rm packed}(y),\beta_{\rm shared}))/\lambda(W)$.
Its domain fits $TD+\kappa$. Packed identities hold in the field ideal before reweighting; shared blocks vanish after multiplying by $W$. Apply the earlier reduction lemmas. Thus, if blocks with $T_a>\tau$ have certified shared cost $\kappa_\tau$, another sufficient bound is $\min_\tau[\tau D+\kappa_\tau]$.
\end{proof}

\section{Limits of pointwise substitutions}
\begin{lemma}[The independent-input degree barrier]\label{lem:pointwisebarrier}
For inputs $g_j=y_j$ on $\F_p^R$, suppose coefficients $\beta_{u,j}(y)$ of degree at most $L$ make all companions vanish at every assignment. Then
\[
 \boxed{L\ge\left\lceil\frac{(p-1)R}{h}\right\rceil-1.}
\]
For independent Boolean inputs over any $\F_p$, the same argument gives $h(L+1)\ge R$.
\end{lemma}
\begin{proof}
Put $P(y)=\prod_u(1-\sum_j\beta_{u,j}y_j)$. The equations $y_jP=0$ force $P(0)=1$ and $P(y)=0$ for all nonzero inputs. Its unique reduced representative on $\F_p^R$ is $\prod_j(1-y_j^{p-1})$, of degree $(p-1)R$. Reduction cannot increase degree, whereas $\deg P\le h(L+1)$. On the Boolean cube the unique multilinear representative is $\prod_j(1-y_j)$, of degree $R$.
\end{proof}

\begin{lemma}[Support on an empty row has full column degree]\label{lem:emptyrowsupport}
In the Boolean column-exclusive quotient $A=k[x]/J$, a nonzero function supported only on assignments where row $i$ is empty has minimum polynomial degree exactly $n$.
\end{lemma}
\begin{proof}
For each column use the basis $1,e_{j,a}$ for all column states $a\ne i$, including the empty state $a=0$. Every $e_{j,a}$ is affine-linear and vanishes at state $i$; conversely $x_{ij}=1-\sum_{a\ne i}e_{j,a}$. Thus the change of basis preserves the degree filtration. Tensoring these column bases gives a unique expansion whose terms select a set of columns.

Set a column $j$ to state $i$. Terms using a nonconstant factor there vanish; all other terms remain. Since the function is zero under this restriction, uniqueness forces every coefficient of a term omitting column $j$ to be zero. Repeat for all columns. Every surviving nonzero term uses all $n$ columns. Such terms have degree $n$, and no lower-degree expansion is possible.
\end{proof}

\begin{theorem}[A barrier to weighted packing modulo column axioms alone]\label{thm:columnbarrier}
For each PHP row $i$ introduce a block with inputs $x_{i1},\ldots,x_{in}$ and common factor
\[
 P_i=\prod_{u=1}^h(1-\sum_jr_{i,u,j}x_{ij}).
\]
Suppose a polynomial $W$ nonzero in $A=k[x]/J$ and substitutions $\beta$ of maximum degree $L$ satisfy
$Wx_{ij}P_i(x,\beta)=0$ in $A$ for all $i,j$. Then
\[
 \deg W+h(L+1)\ge n.
\]
For $D\ge2h+1$ and $T=\max\{1,L\}$, the previous advertised budget obeys
\[
 \boxed{\deg W+TD\ge n+D-2h\ge n+1.}
\]
\end{theorem}
\begin{proof}
Let $Q_i=\prod_j(1-x_{ij})$. The annihilation assumptions imply
$WP_i=Q_iWP_i$, and $Q_iP_i=Q_i$ because $Q_ix_{ij}=0$. Hence
$WP_i=WQ_i$ in $A$. Choose a column-exclusive assignment where $W\ne0$. Since there are $n+1$ rows and only $n$ columns, some row $i$ is empty there; then $WQ_i\ne0$. By \cref{lem:emptyrowsupport},
$n\le\deg(WP_i)\le\deg W+h(L+1)$.
For $L\ge1$, add $DL-h(L+1)$ to this bound and use $D\ge2h+1$; the minimum occurs at $L=1$, giving $n+D-2h$. For $L=0$, the stronger bound $\deg W+D\ge n+D-h$ holds.
\end{proof}

\begin{corollary}[The obstruction is representational, not intrinsic]\label{cor:rowfree}
The same family has a degree-preserving lifting when the row equations are used. Set $r_{i,1,j}=1$ for every $j$ and all other factors' coefficients to zero. Then
\[
 E_{i,j}(x,\beta)=-x_{ij}(\rho_i-1).
\]
Every degree-$D$ base design lifts by constant substitution.
\end{corollary}
\begin{proof}
The selected factor is $1-\rho_i$. Its product with $x_{ij}$ has a degree-two base-axiom certificate, below the companion's original degree $2h+1$. Every permitted multiple therefore lies within the original degree budget. This uses the row axiom, which was deliberately excluded from $J$ in \cref{thm:columnbarrier}.
\end{proof}

\begin{remark}
The lower bound in \cref{thm:columnbarrier} is sharp as a bound on that pointwise packing architecture. For $q=n-2h\ge0$, use $W=\prod_{j=1}^{q}x_{jj}$. On selected rows use $1-x_{ii}$. On each other row pair the remaining $2h$ columns, using coefficients $1$ and $1-x_{ia}$ on a pair $(a,b)$; their factor is $(1-x_{ia})(1-x_{ib})$. This gives $\deg W=n-2h$ and $T=1$. It does not promise $\lambda(W)\ne0$ for a prescribed design; it establishes sharpness of the algebraic degree constraint only.
\end{remark}

\input{05_baseaware.tex}
\input{06_moments.tex}
\chapter{Why PC closure matters, and additive elimination}\label{ch:elimination}
\chapterintro{A long chain of locally simple extension blocks can defeat generic batching from ordinary designs. Polynomial-calculus closure excludes that particular obstruction and leads to an additive proof-elimination theorem. Unlike the preceding moment-preserving constructions, this argument seeks existence of a suitable joint functional, not preservation of a prescribed one.}

\section{An explicit obstruction to batching ordinary designs}
Let $G$ be a DAG with vertices $1,\ldots,N$ in topological order, a unique sink $N$, and indegree at most two. Put
\[
 b_i=1-x_i,\quad f_i=b_i\prod_{j\in\operatorname{pred}(i)}x_j,
 \quad\calF_G=\{f_i\}_i\cup\{x_N\}\cup\{x_i^2-x_i\}_i.
\]
For every prefix $[k]$ add one block with inputs $b_1,\ldots,b_k$, accuracy $h$, and notation
\[
 A_k=\prod_{u=1}^h(1-\sum_{j\le k}r_{k,u,j}b_j),\qquad
 E_{k,j}=b_jA_k,\qquad A_0=1.
\]
All blocks are in one level.

\begin{lemma}[A degree-$(4h+2)$ prefix certificate]\label{lem:prefixcertificate}
The augmented pebbling system has an exact NS refutation of degree at most $4h+2$, over every field. Each certificate summand uses variables from at most two consecutive blocks. Boolean and extension field equations are not needed in the certificate.
\end{lemma}
\begin{proof}
Define
\[
 U_{k,j}=\sum_{u=1}^h r_{k,u,j}
 \prod_{v<u}(1-\sum_{\ell\le k}r_{k,v,\ell}b_\ell).
\]
Product telescoping gives $1-A_k=\sum_{j\le k}U_{k,j}b_j$ and $\deg U_{k,j}\le2h-1$. Thus
\begin{align*}
 A_{k-1}-A_k={}&\sum_{j<k}U_{k,j}E_{k-1,j}
 +U_{k,k}b_kA_{k-1}
 -\sum_{j<k}U_{k-1,j}E_{k,j}.
\end{align*}
Order the at most two predecessors of $k$ as $p_1,\ldots,p_t$. The local identity
\[
 b_kA_{k-1}=A_{k-1}f_k+
 b_k\sum_{\ell=1}^t\left(\prod_{v<\ell}x_{p_v}\right)E_{k-1,p_\ell}
\]
follows by expanding $1-\prod_{\ell=1}^tx_{p_\ell}$. Substitution expresses every consecutive difference as an axiom combination. Its largest summand degree is
$(2h-1)+2h+3=4h+2$.

The differences telescope to $1-A_N$, while
$A_N=A_Nx_N+E_{N,N}$. Therefore
\[
 1=\sum_{k=1}^N(A_{k-1}-A_k)+A_Nx_N+E_{N,N}
\]
is the required certificate after substituting the local expansions. The formulas show the adjacent-block support and do not use field reduction.
\end{proof}

\begin{imported}[Pebbling degree input]\label{imp:pebbling}
For the pebbling encoding above, NS degree equals reversible pebbling space, over every field \cite[Theorem 3.1]{Pebbling}. The Carlson--Savage constructions provide bounded-indegree single-sink families with polynomially growing pebbling price; the conversation used a family with $N=\Theta(r^3)$ and price $\Omega(r)$. These are published combinatorial inputs, not results proved in this document.
\end{imported}

\begin{corollary}[Plain high-degree designs do not justify universal batching]\label{cor:genericbatchfalse}
There is no theorem valid for every Boolean base system asserting that a base design through $(D+h+\log M)^C$, for a fixed constant $C$, always yields an augmented degree-$D$ design for one-level ENS systems with $M$ companions.
\end{corollary}
\begin{proof}
Use a family from Imported Input~\ref{imp:pebbling}. Its base has ordinary designs through polynomially growing degree below its NS threshold. The prefix construction has $M=N(N+1)/2$ companions and degree-$D$ refutations for $D=4h+2$. Take $h=\Theta(\log N)$, so the putative sufficient starting degree is only polylogarithmic and is eventually below the base NS threshold. The asserted augmented design would contradict \cref{lem:prefixcertificate}. Each certificate summand still uses at most two adjacent blocks, so sparse local interaction alone does not rescue that theorem.
\end{proof}

\begin{lemma}[The same pebbling base has PC degree at most three]\label{lem:pebblingpc}
The base system $\calF_G$ has a PC refutation in degree at most three.
\end{lemma}
\begin{proof}
For sources, $b_i=f_i$ is an axiom. In topological order use
\[
 b_k=f_k+b_k\sum_{\ell=1}^t b_{p_\ell}
              \prod_{v<\ell}x_{p_v}.
\]
The predecessor $b_{p_\ell}$ are already derived. Multiplying their final lines as displayed costs degree at most three since $t\le2$. This derives every $b_k$, and $b_N+x_N=1$ finishes. Thus the high NS degree of the example does not imply high PC degree.
\end{proof}

\section{A stronger input functional}
\begin{lemma}[PC-based supported lifting]\label{lem:pcbased}
For a block of maximum input degree $\delta$ over $\F_p$, put
\[
 \Psi_{p,\delta}(D)=\max\{D,\ D-h+(p-1)\delta,
                      \ p(D-h(\delta+1))\}.
\]
A normalized functional annihilating $\Cc{\Psi_{p,\delta}(D)}(\calF)$ gives an ordinary degree-$D$ augmented design, preserving old moments. The output is not asserted to annihilate all augmented PC consequences.
\end{lemma}
\begin{proof}
Repeat the supported construction of \cref{lem:supported-p}, but require components on degree $N_j=D-h+(p-1)\delta_j$ to annihilate $\Cc{N_j}(\calF)$ instead of $\mathcal I_{N_j}$. In the dual kernel condition, every $Gg_j$ is already derivable through $N_j$. Write $a=h(\delta+1)$ and $G=\sum_iq_ig_i$ with $\deg G\le D-a$. By \cref{lem:reuse},
\[
 G^p=\sum_iq_iG^{p-2}(Gg_i)
\]
can be derived through $\max\{\max_iN_i,p(D-a)\}$: multiply the derived polynomials, not their flattened certificates. Field reduction derives $G^p-G$ within $p(D-a)$. Thus the input functional annihilates every dual kernel element, giving component consistency. The full cofactor verification is unchanged from \cref{lem:supported-p}, and yields an ordinary design. Nothing in that verification establishes closure under reuse of augmented derived polynomials, so the stronger output property is not inferred.
\end{proof}

For $p=2$, this gives $\Psi_\delta(D)=\max\{D,D-h+\delta,2D-2h(\delta+1)\}$. It is degree-preserving when $\delta\le h$ and $D\le2h(\delta+1)$. This intermediate lemma exposed the useful closure distinction; the following proof transformation handles composition directly.

\section{Eliminate one arbitrary block}
\begin{theorem}[One-block additive elimination]\label{thm:one-elimination}
Let $\calG\subseteq\F_p[y]$ contain Boolean or field equations for every old variable. Suppose an accuracy-$h$ block has fresh variables absent from $\calG$, and maximum input degree $\delta$. A degree-$d$ PC refutation of $\calG\cup\calE\cup\calR$ yields a PC refutation of $\calG$ of degree at most
\[
 \boxed{d+(p-1)\delta.}
\]
There is no fan-in factor.
\end{theorem}
\begin{proof}
Keep the original refutation $\pi$. Fix an input $g_j$ and $\alpha\ne0$. Set $r_{1,j}=\alpha^{-1}$ and every other new variable zero. Multiply each specialized proof line by $\chi_\alpha(g_j)$. An extension axiom maps to
\[
 \chi_\alpha(g_j)g_i(1-\alpha^{-1}g_j)
 =\alpha^{-1}g_i(g_j^p-g_j),
\]
a field-ideal polynomial. Its actual degree is bounded by the degree of the transformed original axiom line, hence by $d+(p-1)\delta$. \Cref{lem:fieldreduction} supplies a derivation in that budget. Old axioms become old-axiom multiples, and new field equations vanish.

By \cref{lem:substitution}, all weighted specialized inference lines can be derived within the same degree. Since the last line is one, this gives a derivation of $\chi_\alpha(g_j)$ from $\calG$. Repeat over $\alpha\ne0$ and use \eqref{eq:selectorlinear} to derive $g_j$. Do this for every input. Multiple derivations increase length, not the maximum degree.

Now specialize the original $\pi$ with all new variables zero. Every extension axiom becomes $g_i$, already derived. Reuse these final polynomials as the required starting lines. The specialized inference sequence has degree at most $d$, and the inserted input derivations have degree at most $d+(p-1)\delta$. This refutes $\calG$ within the claimed budget.
\end{proof}

\begin{theorem}[Additive elimination through all levels]\label{thm:additive}
For a leveled system with block input degrees $\delta_1,\ldots,\delta_s$,
\[
 \boxed{\deg_{\PC}(\calF\cup\calE\cup\calR)
 \ge\deg_{\PC}(\calF)-(p-1)\sum_{a=1}^s\delta_a.}
\]
Equivalently, a degree-$D$ augmented PC refutation yields a base refutation of degree at most $D+(p-1)\sum_a\delta_a$.
\end{theorem}
\begin{proof}
Eliminate blocks in reverse level order. Later blocks have already disappeared, so the fresh variables of the block being removed occur only in that block and its own field equations. Take all remaining axioms as the old system and apply \cref{thm:one-elimination}. Its output is an actual PC refutation, exactly the object required for the next step. Add the degree increments.
\end{proof}

\begin{corollary}[Existence of a PC-annihilating joint functional]\label{cor:pcdesign}
If the base has no PC refutation through degree $D+(p-1)\sum_a\delta_a$, then a normalized functional on joint polynomials of degree at most $D$ annihilates $\Cc D(\calF\cup\calE\cup\calR)$. In particular it is an ordinary augmented design.
\end{corollary}
\begin{proof}
Otherwise \cref{thm:additive} would yield a forbidden base refutation. Apply \cref{lem:duality} to the augmented PC consequence space. This is existence of a suitable functional, not extension of every preassigned base functional.
\end{proof}

\begin{lemma}[Equal-span batches]\label{lem:equalspan}
Suppose several same-level blocks have identical polynomial input span $V$, with a basis of degree at most $\delta_V$. They can be eliminated simultaneously at degree cost $(p-1)\delta_V$. If $1\in V$, scalar specialization eliminates the whole batch at zero degree cost.
\end{lemma}
\begin{proof}
For each basis polynomial $f\in V$ and each block, express $f$ as a constant linear combination of that block's inputs. For a fixed $\alpha\ne0$, put the scaled coefficients into the first extension factor of every block, making the selected factor $1-\alpha^{-1}f$ everywhere. Weight the specialized original proof by $\chi_\alpha(f)$. All selected extension axioms become field consequences simultaneously. Derive the basis inputs as in \cref{thm:one-elimination}, then derive every block input by constant linear combination. Zero-specialize the original proof and reuse them. If $1\in V$, choose coefficients making the first factor identically zero in every block, with no weight.
\end{proof}

\begin{remark}
For $p=2$, linear inputs, $h=16$, $D=49$, ten unrelated blocks cost at most $59$ by additive elimination; ten blocks of a common span cost $50$. The earlier moment-preserving recurrence gave $48+2^{10}=1072$ for ten sequential blocks at this boundary. These are different sufficient bounds, not statements that the higher cost is necessary.
\end{remark}

\chapter{Nested spans and core--residual batches}\label{ch:batching}
\chapterintro{The original refutation can be replayed while previously learned inputs are reused. This removes chain length from the degree cost. Polynomial factor packing then allows the full spaces to be incomparable, provided suitable cores are nested and the remaining directions are inexpensive.}

\section{A chain costs one block}
\begin{theorem}[Nested-span batch elimination]\label{thm:nested}
Let same-level blocks, fresh for the old system $\calG$, have input spans
\[
 V_1\subseteq V_2\subseteq\cdots\subseteq V_s.
\]
If all inputs have degree at most $\delta$, a degree-$D$ augmented PC refutation yields a refutation of $\calG$ of degree at most
\[
 \boxed{B=D+(p-1)\delta.}
\]
The bound is independent of the length of the chain and all fan-ins.
\end{theorem}
\begin{proof}
Fix the original degree-$D$ refutation $\pi$ and never replace it by an inflated proof. Process blocks in increasing order, maintaining derivations of the earlier inputs within degree $B$.

Choose an input $f=g_{a,j}$ of the current block. For every later block $b\ge a$, nesting gives constants $c_{b,i}$ with $f=\sum_i c_{b,i}g_{b,i}$. For a fixed $\alpha\ne0$, zero-specialize all earlier blocks; in every current or later block set $r_{b,1,i}=\alpha^{-1}c_{b,i}$ and other factor coefficients zero. Weight each specialized line of the original $\pi$ by $\chi_\alpha(f)$.

A current or later extension axiom becomes
$\alpha^{-1}g_{b,i}(f^p-f)$, an old-field consequence within the weighted-line degree ceiling $B$. An earlier axiom becomes $\chi_\alpha(f)g_{b,i}$, where $g_{b,i}$ is already derived. Reuse its final polynomial and multiply it by the selector; the product degree is at most $B$ for every axiom line occurring in $\pi$. The earlier derivation itself is not multiplied. Old axioms, new field axioms, and inferences are handled by \cref{lem:substitution,lem:fieldreduction}.

The transformed conclusion is $\chi_\alpha(f)$. Combining these derivations over $\alpha\ne0$ derives $f$ within $B$. Repeat for the current inputs and move on. The same original proof and the same ceiling $B$ are used at every learning stage.

Finally zero-specialize all new variables in $\pi$. Every extension axiom is now an already derived input. Reuse the derivations to get a base refutation within $B$. Unused axioms need no substitution certificate; a block absent from the used extension axioms can simply be zero-specialized.
\end{proof}

\begin{lemma}[Nested modulo supplied base consequences]\label{lem:nestedquotient}
Let $H\subseteq k[y]_{\le\delta}$ be a vector space whose elements have supplied PC derivations from $\calG$ within degree $c$. Suppose
\[
 (V_1+H)/H\subseteq\cdots\subseteq(V_s+H)/H.
\]
The batch can be eliminated within degree
\[
 \boxed{B_H=\max\{c,D+(p-1)\delta,(p+1)\delta\}.}
\]
If $c\le D$ and $D\ge2\delta$, this is $D+(p-1)\delta$.
\end{lemma}
\begin{proof}
A selected input now has $f=\sum_i c_{b,i}g_{b,i}+u_b$ with $u_b\in H$. The same first-factor specialization gives
\[
 \chi_\alpha(f)E_{b,i}(y,\beta)
 =\alpha^{-1}g_{b,i}(f^p-f)
  +\alpha^{-1}\chi_\alpha(f)g_{b,i}u_b.
\]
The first term is a field consequence, the second a multiple of an already derived $u_b$. Both have degree at most $(p+1)\delta$. Reuse their derivations within $B_H$; the rest of \cref{thm:nested} is unchanged. The quotient here is a vector-space quotient backed by low-degree derivations, not the full ideal of the unsatisfiable base.
\end{proof}

\begin{corollary}[Global chain charge]\label{cor:chaincharge}
Partition every level into nested chains $C$, with maximum input degree $\delta_C$. Eliminating in reverse level order yields
\[
 \boxed{D_{\rm base}\le D+(p-1)\sum_C\delta_C.}
\]
The quotient version uses the extra ceilings in \cref{lem:nestedquotient}. Literal equality of all spaces is the special case of a constant chain.
\end{corollary}
\begin{proof}
After later levels have been removed, a chosen chain is fresh for all other remaining axioms. Apply \cref{thm:nested}, treating those axioms as old. Scalar substitutions do not increase the degrees of surviving blocks. Add the increments.
\end{proof}

For example, $1{,}000$ strictly nested linear spaces over $\F_2$ at target degree $49$ cost degree $50$, rather than the per-block budget $1049$. This is a degree transformation; its proof length is not bounded here.

\section{Nested cores with incomparable full spaces}
\begin{theorem}[Core--residual batch elimination]\label{thm:cores}
For same-level input spaces $V_a$, suppose there are nested cores
\[
 U_1\subseteq\cdots\subseteq U_s,\qquad U_a\subseteq V_a,
\]
each generated in degree at most $\gamma$. Choose representatives $v_{a,1},\ldots,v_{a,r_a}\in V_a$ for a basis of $V_a/U_a$, with degrees at most $\eta_a$. Put
\[
 T_a=\max\left\{1,\eta_a
 \left(\left\lceil\frac{(p-1)r_a}{h_a}\right\rceil-1\right)_+\right\},
 \quad T=\max_aT_a.
\]
When $r_a=0$ take $\eta_a=0$. A degree-$D$ augmented PC refutation yields a base refutation of degree at most
\[
 \boxed{B=TD+(p-1)\gamma.}
\]
The full spaces $V_a$ need not be nested.
\end{theorem}
\begin{proof}
For one block, decompose
\[
 g_{a,i}=u_{a,i}+\sum_\nu c_{i,\nu}v_{a,\nu},\qquad u_{a,i}\in U_a,
\]
and define $Z_a=\prod_\nu(1-v_{a,\nu}^{p-1})$.
By the packed coefficient construction, using the residual representatives rather than a basis of all inputs, there are substitutions $\beta_a(y)$ of degree at most $T_a$ making the common extension factor exactly $Z_a$. The representatives are constant linear combinations of the original tuple, so these are legitimate coefficients in its original input coordinates. Then
\begin{equation}\label{eq:coreabsorb}
 E_{a,i}(y,\beta_a)=g_{a,i}Z_a
 =u_{a,i}Z_a+(g_{a,i}-u_{a,i})Z_a,
\end{equation}
and the last term belongs to $\Jfld$.

Keep the original refutation $\pi$. Process cores in increasing order. Assume the earlier cores and earlier specialized extension axioms have been derived within $B$. Select a core generator $f\in U_a$ of degree at most $\gamma$. For every $b\ge a$, $f\in V_b$, so choose constant coordinates of $f$ in that tuple. Substitute the residual polynomial coefficients in earlier blocks, and in every current or later block use the scalar first-factor coefficients making $1-\alpha^{-1}f$. Weight the image of every original line by $\chi_\alpha(f)$.

All substituted variables have degree at most $T$. Thus every transformed line has degree at most $TD+(p-1)\gamma$. Current and later extension images are old-field consequences by \eqref{eq:selectorproduct}. Earlier images are already derived and their final polynomials can be multiplied by the selector. Old axioms and field equations map to bounded-degree consequences. For an original multiplication step, the transformed predecessor has degree at most $T(D-1)+(p-1)\gamma$; multiplication by an image of degree at most $T$ stays within $B$. This explicitly validates the use of polynomial, rather than scalar, substitutions.

As before, the transformed conclusion derives $\chi_\alpha(f)$, and combining the nonzero field values derives $f$ within $B$. Learn generators of the core this way.

To finish the block, use \eqref{eq:coreabsorb}. The polynomial $u_{a,i}$ is now derived and has degree at most $\gamma$. For an extension axiom actually used in $\pi$, $\deg E_{a,i}\le D$, hence $\deg(g_{a,i}Z_a)\le TD$. If $Z_a=0$ the image is zero. Otherwise the ordinary polynomial ring is a domain, so $\deg Z_a\le TD$. Therefore both $u_{a,i}Z_a$ and $(g_{a,i}-u_{a,i})Z_a$ have degree at most $TD+\gamma\le B$. The former is derived by reuse, the latter by bounded-degree field reduction. This derives the specialized extension image within the same ceiling.

Continue through all cores and finally substitute all residual coefficients into the original $\pi$. Every extension image has a supplied derivation and every inference is simulated within $B$, yielding the extension-free refutation.
\end{proof}

\begin{corollary}[Cheap incomparable families]\label{cor:cheapantichain}
For affine-linear inputs, residual rank
$r_a\le\lfloor2h_a/(p-1)\rfloor$ gives $T=1$. Thus nonzero affine cores cost only $p-1$ degrees per batch. For arbitrary nonlinear residual inputs, the stronger rank condition $r_a\le\lfloor h_a/(p-1)\rfloor$ also gives $T=1$, and the cost is $(p-1)\gamma$ regardless of their degrees.
\end{corollary}
\begin{proof}
In the affine case $\eta_a\le1$ and the ceiling in \cref{thm:cores} is at most two. In the nonlinear case it is at most one, so every residual substitution is scalar. Apply the theorem.
\end{proof}

An example is $V_a=U+P_a$ with the same arbitrarily large affine core $U$ and different residual subspaces of dimension at most $2h$ over $\F_2$. The full spaces can form an arbitrarily large antichain, yet the entire batch costs at most $D+1$. At $h=16,D=49$, residual rank $32$ is allowed while full-group interactions are already active.

\begin{lemma}[Affine cores modulo base consequences]\label{lem:corequotient}
Let $H\subseteq k[y]_{\le1}$ have supplied derivations within degree $c$. Form affine input spaces modulo $H$, and suppose they have nested cores and residual representatives as above. For $D\ge2$ a sufficient budget is
\[
 \boxed{\max\{c,TD+p-1\}.}
\]
\end{lemma}
\begin{proof}
Lift each quotient core generator $f$ to an affine representative. In a later block write $f=\sum_i c_i g_i+u$ with $u\in H$. The weighted image of a selected companion is the sum of
$\alpha^{-1}g_i(f^p-f)$ and $\alpha^{-1}\chi_\alpha(f)g_i u$.
Their degrees are at most $p+1\le TD+p-1$; the second uses the supplied derivation of $u$. For residual absorption, the discrepancy between an input and its residual combination belongs to the learned core plus $H$, so it is already derivable within the stated ceiling. The rest is \cref{thm:cores}.
\end{proof}

For PHP one may take $H=\Span\{\rho_i-1\}$, whose generators are degree-one axioms. Larger supplied consequence spaces are allowed only when their derivations are genuinely available within the asserted degree.

\begin{lemma}[Maximal cores for a fixed order]\label{lem:maxcores}
For an order $a_1,\ldots,a_t$, the maximal possible nested cores are
\[
 U_{a_i}=\bigcap_{j=i}^tV_{a_j}.
\]
Every other nested core choice satisfying $U_{a_i}\subseteq V_{a_i}$ is contained in these suffix intersections.
\end{lemma}
\begin{proof}
If $U_{a_i}$ is nested, it lies in $U_{a_j}\subseteq V_{a_j}$ for every $j\ge i$, so it is contained in the intersection. The suffix intersections themselves form an increasing chain and lie in the corresponding spaces.
\end{proof}

\begin{remark}
A batch transforms a degree budget as $d\mapsto T_{\calB}d+(p-1)\gamma_{\calB}$. For $T_{\calB}=1$, charges add; otherwise they compose as affine functions and can multiply earlier budgets. A constant number of levels with polylogarithmic $T$ and core degree would be affordable. No such uniform decomposition was proved. The next chapter determines the exact optimum of this criterion for affine data and shows it can exceed the PHP budget.
\end{remark}

\chapter{Exact decomposition and a sharp limitation}\label{ch:decomposition}
\chapterintro{For affine inputs, both the ordering and the optimal partition cost of the current core--residual rule can be determined exactly. Nevertheless, a polynomial-size admissible family has optimum $D+(p-1)n$, exceeding the available budget. This is a limitation of the rule, not of every elimination method.}

\section{Ordering is solved}
Assume one level, common accuracy $h$, and affine input spaces $V_1,\ldots,V_s$ over $\F_p$. Supplied affine consequences can first be quotiented out. Zero spaces are discarded; spaces containing a nonzero constant are handled by the separate free-elimination rule.

For a family $S$ and an order $a_1,\ldots,a_t$, define its residual radius
\[
 R(a_1,\ldots,a_t)=\max_i\left[
 \dim V_{a_i}-\dim\left(\bigcap_{j=i}^tV_{a_j}\right)\right].
\]
Let $R_*(S)$ be the minimum over orders.

\begin{lemma}[Optimal order and exact obstruction formula]\label{lem:ordering}
Ordering by nondecreasing dimension minimizes the residual radius. Moreover,
\[
 \boxed{R_*(S)=\max_{\varnothing\ne A\subseteq S}
 \left[\min_{a\in A}\dim V_a-\dim\left(\bigcap_{a\in A}V_a\right)\right].}
\]
\end{lemma}
\begin{proof}
For the exchange argument, consider adjacent spaces $A,B$ with $\dim A\ge\dim B$, and let $W$ be the intersection of the spaces after them. Before swapping, the residual of $A$ is $\dim A-\dim(A\cap B\cap W)$. After swapping the two affected residuals are
\[
 \dim B-\dim(A\cap B\cap W),\qquad
 \dim A-\dim(A\cap W),
\]
both at most that previous residual. All other terms are unchanged. Removing inversions therefore yields a nondecreasing-dimension optimum.

For the formula's lower bound, in any order take the first member of a nonempty subfamily $A$. Its suffix contains all of $A$, so its residual is at least the bracketed quantity. For the upper bound, repeatedly remove a smallest-dimensional remaining space. At each step its residual is exactly the bracketed quantity for the remaining family, hence at most the maximum over all subfamilies.
\end{proof}

\section{Optimize the whole level}
Put
\[
 T(r)=\max\left\{1,\left\lceil\frac{(p-1)r}{h}\right\rceil-1\right\}.
\]
\begin{lemma}[The optimal certified batch cost]\label{lem:batchcost}
Within the affine core--residual criterion, the best budget for a nonempty batch $S$ acting on degree $d$ is
\[
 \boxed{C_S(d)=\min\{T(R_*(S))d+(p-1),\
                       T(\max_{a\in S}\dim V_a)d\}.}
\]
The alternatives are maximal nonzero cores and all-zero cores. This optimizes the stated sufficient bound, not all possible proof transformations.
\end{lemma}
\begin{proof}
Every nonzero affine core has generator degree one and therefore charge $p-1$. Maximal suffix cores minimize every residual for a fixed order, and \cref{lem:ordering} minimizes their maximum. Shrinking some but not all cores cannot lower that charge and cannot improve the residual maximum. If all cores are zero, there is no core charge but each residual rank is the whole input rank. These give the two displayed alternatives.
\end{proof}

\begin{lemma}[Exact subset optimization]\label{lem:subsetdp}
The best degree bound supplied by this rule over all partitions and batch orders of a family $S$ is given by
\[
 \boxed{B(\varnothing)=D,\qquad
 B(S)=\min_{\varnothing\ne A\subseteq S}C_A(B(S\setminus A)).}
\]
This is an exact finite algorithm, generally exponential in the number of blocks.
\end{lemma}
\begin{proof}
Choose the last eliminated batch $A$. Every $C_A$ is nondecreasing, so among eliminations of the preceding blocks only the smallest resulting degree can improve the final answer. Induction on $|S|$ proves the recurrence. Conversely, every choice in the recurrence corresponds to a partition and batch order, establishing attainability of its certified value.
\end{proof}

\section{An admissible family with unaffordable optimum}
Let $m=n+1$ and
\[
 H_0=\Span_{\F_p}\{\rho_i-1:i\in[n+1]\},\qquad
 W=\Span_{\F_p}\{x_{ij}:i\in[n+1],j<n\}.
\]
\begin{lemma}[A large affine coordinate space survives the row quotient]\label{lem:coordinatequotient}
The $n^2-1$ displayed coordinate variables, together with one, are linearly independent modulo $H_0$.
\end{lemma}
\begin{proof}
If $c+w=\sum_i a_i(\rho_i-1)$ with $w\in W$, compare coefficients of $x_{i,n}$: every $a_i=0$. Then $c+w=0$ forces $c=0$ and every coefficient of $w$ to be zero.
\end{proof}

\begin{lemma}[A polynomial-size spread of input spaces]\label{lem:spread}
Set
\[
 d=\lceil\log_p n\rceil,\qquad
 q=\left\lfloor\frac{n^2-1}{2d}\right\rfloor,\qquad r=qd,
\]
for $n$ large enough that $d,q\ge1$. There exist $n$ admissible same-level input spaces of dimension $r$ in the original PHP variables, pairwise intersecting only at zero, even modulo $H_0$.
\end{lemma}
\begin{proof}
Let $K=\F_{p^d}$ and choose $n$ distinct scalars $\alpha\in K$. Embed the $2r$-dimensional $\F_p$-space $K^q\oplus K^q$ into $W$. Define
\[
 V_\alpha=\{(u,\alpha u):u\in K^q\}.
\]
Each has $\F_p$-dimension $r$. If $(u,\alpha u)=(v,\beta v)$ and $\alpha\ne\beta$, then $u=v$ and $(\alpha-\beta)u=0$, forcing $u=0$. \Cref{lem:coordinatequotient} preserves this intersection property after the row quotient.

Choose a basis of each space as its homogeneous linear input tuple and introduce a fresh accuracy-$h$ block with all its companions. This is valid ENS syntax. There are $n$ blocks, $nr=O(n^3)$ companions, and $hnr$ new variables. Each companion has original degree $2h+1$.
\end{proof}

\begin{theorem}[The exact best decomposition bound for the spread]\label{thm:spreadcost}
For the family in \cref{lem:spread}, the best degree bound supplied by the current affine criterion is
\[
 \boxed{B_{\rm criterion}^{\rm opt}
       =\min\{D+(p-1)n,T(r)D\}.}
\]
For fixed $p$, $h=\Theta(\log n)$, and $2h+1\le D=\operatorname{polylog}n$, this equals $D+(p-1)n>n/2$ for all sufficiently large $n$.
\end{theorem}
\begin{proof}
In any batch of at least two spaces, the first core is contained in the intersection with a later space and is therefore zero. Its residual already has rank $r$, so $R_*=r$ for every such batch. The best multi-block cost is $T(r)d$. A singleton offers $\min\{d+p-1,T(r)d\}$.

If every singleton uses a full core, the cost is $D+(p-1)n$. Any multi-block batch, or any all-zero-core singleton, incurs a multiplicative step $T(r)$ and therefore costs at least $T(r)D$ overall; all other steps are nondecreasing. One all-zero-core batch attains this latter bound. This proves the exact minimum over all allowed decompositions.

For fixed $p$, $r=\Theta(n^2)$ and $T(r)=\Theta(n^2/\log n)$. Since $D\ge2h+1=\Theta(\log n)$, $T(r)D=\Omega(n^2)$, while $D+(p-1)n=\Theta(n)$. The first branch of the minimum eventually wins and exceeds $n/2$.
\end{proof}

\begin{example}[A numerical instance of the criterion]
At $p=2,n=256,h=8,D=25$, the construction has $r=32760$ and $T(r)=4094$. The two optimized certified costs are $281$ and $102350$, while the available target budget is $128$. This is an illustration of the degree-bound criterion, not a claimed ENS refutation or a fully matched Frege-simulation instance.
\end{example}

\section{Low-degree affine consequences do not improve this quotient}
\begin{theorem}[Affine-consequence rigidity, conditional on the PHP degree input]\label{thm:rigidity}
Using Imported Input~\ref{imp:php} for the base and the residual PHP systems, for
$1\le b\le\lfloor(n-2)/2\rfloor$,
\[
 \boxed{\Cc b(\Fphp)\cap\F_p[x]_{\le1}=H_0.}
\]
\end{theorem}
\begin{proof}
Let $f=a_0+\sum_{i,j}a_{ij}x_{ij}$ have a degree-$b$ PC derivation. Fix distinct rows $i,i'$ and columns $j,k$. Restrict the derivation to the partial matching $i\mapsto j,i'\mapsto k$, and to the matching $i\mapsto k,i'\mapsto j$. Both leave exactly the same labeled residual PHP system. The two restricted affine polynomials differ by the constant
\[
 a_{ij}+a_{i'k}-a_{ik}-a_{i'j}.
\]
A nonzero difference would give a degree-$b$ residual refutation by subtraction and scalar division, forbidden by the residual lower bound. Thus all coefficient rectangles vanish, which implies $a_{ij}=u_i+v_j$ (fix one row and column to construct $u_i,v_j$).

Subtract a combination of row axioms. It remains to consider a derived polynomial $f'=c+\sum_jv_j\sigma_j$, where $\sigma_j=\sum_i x_{ij}$. If $v_j\ne v_k$, permute columns $j,k$ in its derivation and subtract. Dividing by $v_j-v_k$ derives $\sigma_j-\sigma_k$ within degree $b$.

Match one pigeon to column $j$. This yields $1-\sigma'_k$ in the residual system. Permuting its remaining columns derives $1-\sigma'_\ell$ for all $n-1$ residual holes. Sum these polynomials and the $n$ residual row equations:
\[
 \sum_{\ell=1}^{n-1}(1-\sigma'_\ell)
 +\sum_{i=1}^{n}(\rho'_i-1)=-1.
\]
This is another forbidden residual refutation. Therefore all $v_j$ agree. Row equations reduce $f'$ to a constant, which must be zero by the base lower bound. Hence $f\in H_0$. The opposite inclusion consists of base axioms and their linear combinations.
\end{proof}

\begin{remark}
Thus adding all base-PC affine consequences within the displayed range creates no new affine identifications beyond the row equations. In particular, polylogarithmic-degree affine consequences cannot improve the spread decomposition. This does not exclude nonlinear reductions, relations using remaining extension blocks, or transformations informed by the actual refutation.
\end{remark}

\section{The precise boundary reached}
The decomposition problem for the stated affine criterion is solved, but its universally affordable version is false. The spread construction supplies admissible data, \emph{not} a refutation using those data. It does not show that the blocks occur essentially in the translation of a short Frege proof. Nor does the criterion's optimum lower-bound all possible elimination methods.

The next theorem must therefore be sensitive to the refutation: discard or rewrite irrelevant blocks, exploit the coefficients and cofactors of a translated NS certificate, or use PHP structure outside static affine spaces. A better ordering of the same spaces cannot by itself repair the example; ordering and partition costs have already been optimized.

\input{10_route.tex}
\appendix
\titleformat{\chapter}[display]{\normalfont\huge\bfseries\color{ink}}{\small\sffamily\bfseries APPENDIX \thechapter}{7pt}{}
\input{11_appendices.tex}
\begin{thebibliography}{9}
\addcontentsline{toc}{chapter}{References}
\bibitem{BIKPRS} S. R. Buss, R. Impagliazzo, J. Kraj\'i\v{c}ek, P. Pudl\'ak, A. A. Razborov, and J. Sgall. \emph{Proof complexity in algebraic systems and bounded depth Frege systems with modular counting}. Computational Complexity 6(3), 256--298, 1996/1997. The simulation used here is cited through Kraj\'i\v{c}ek's Theorem 5.2 below.
\bibitem{Krajicek} J. Kraj\'i\v{c}ek. \emph{Extended Nullstellensatz proof systems}. arXiv:2301.10617v3, 2023. Definition 2.1, Lemma 4.1, and Theorem 5.2. \href{https://arxiv.org/abs/2301.10617}{arXiv record}.
\bibitem{Razborov} A. A. Razborov. \emph{Lower bounds for the polynomial calculus}. Computational Complexity 7, 291--324, 1998. DOI: 10.1007/s000370050013. \href{https://doi.org/10.1007/s000370050013}{Publisher record}.
\bibitem{Pebbling} S. F. de Rezende, O. Meir, J. Nordstr\"om, and R. Robere. \emph{Nullstellensatz Size-Degree Trade-offs from Reversible Pebbling}. arXiv:2001.02481; preliminary version, CCC 2019. Theorem 3.1 and the Carlson--Savage graph discussion. \href{https://arxiv.org/abs/2001.02481}{arXiv record}.
\end{thebibliography}
\end{document}
